\documentclass{article}
\usepackage[paperwidth=16in,paperheight=10.5in,margin=0.42in]{geometry}
\usepackage[T1]{fontenc}\usepackage{lmodern,amsmath,booktabs,tabularx,array}
\usepackage[table]{xcolor}
\renewcommand{\familydefault}{\sfdefault}\pagestyle{empty}\setlength{\parindent}{0pt}
\definecolor{maroon}{HTML}{890B23}\definecolor{pale}{HTML}{F8EDEF}
\definecolor{stripe}{HTML}{F4F5F6}
\newcolumntype{Y}{>{\raggedright\arraybackslash}X}
\setlength{\tabcolsep}{9pt}\renewcommand{\arraystretch}{1.24}
\begin{document}
{\fontsize{30}{35}\selectfont\bfseries\color{maroon} Structural descriptors explored\par}
\vspace{10pt}
{\fontsize{16}{20}\selectfont Historical inventory grouped by the physical quantity being assessed.}\par
\vspace{15pt}
\fontsize{15.5}{20}\selectfont
\rowcolors{2}{stripe}{white}
\begin{tabularx}{\textwidth}{@{}>{\bfseries\raggedright\arraybackslash}p{4.2in}Y>{\raggedleft\arraybackslash}p{1.1in}@{}}
\toprule
\rowcolor{pale} Descriptor family & \textbf{Representative quantities} & \textbf{Distinct definitions}\\
\midrule
Activation geometry and catalytic environment & Attack angles and near-attack fractions; local metal geometry; reactive-group exposure; nearby charged residues. & 10 \\[3pt]
Amine--acyl approach geometry (step 2) & Amine N--acyl C distance; B\"urgi--Dunitz angle; near-attack fraction; accepted-pose count. & 4 \\[3pt]
Reactive-group preorganization & Carboxylate positional spread; attack-direction variability; mean angle; bidentate-contact fraction. & 4 \\[3pt]
Hypothesis-specific geometric descriptors & Oxyanion contacts; substrate exposure and access; amine approach; adenine burial; accepted-pose count. & 9 \\[3pt]
Ligand exposure and pocket engagement & Solvent-accessible surface areas of selected reactive atoms; intermediate--protein separation. & 5 \\[3pt]
Substrate-access tunnel geometry & Tunnel enclosure and length under alternative clearance thresholds. & 6 \\[3pt]
Amine-site electrostatic environment & Nearby acidic oxygens; distance to a basic residue; net charge within alternative probe radii. & 12 \\[3pt]
Electrostatic potential and electric field & Potential and projected field at ATP P$_\alpha$ and the acyl carbonyl; ensemble means and variation. & 8 \\[3pt]
Ionizable-residue and p$K_a$ descriptors & Histidine p$K_a$ proximity to 7; minimum lysine p$K_a$; ionizable-residue counts and ensemble variation. & 12 \\[3pt]
Magnesium coordination geometry & Mg--phosphorus and Mg--oxygen distances; coordination and bridging fractions across distance thresholds. & 10 \\[3pt]
Protein--ligand interface confidence & Directional pairwise ipTM; symmetric and extrema-based summaries across predicted poses. & 7 \\[3pt]
Minimum substrate--ATP distance & Minimum donor oxygen--reactive phosphorus distance in the corresponding prediction ensemble. & 1 \\[3pt]
\midrule
\rowcolor{pale}\textbf{Total} & Distinct, nonconstant definitions in the historical inventory & \textbf{88}\\
\bottomrule\end{tabularx}\par\vspace{14pt}
{\fontsize{11.5}{15}\selectfont
\textbf{Counting rule:} 98 exported fields = 88 distinct nonconstant definitions + 5 duplicate exports + 5 constant or undefined fields.
Counts include formula variants, threshold variants and sampling diagnostics; they are not counts of independent mechanistic hypotheses.
Related quantities can appear in different groups when their definitions or modeling setups differ. Later follow-up formulas overlap this inventory and are not added to these counts.\\[5pt]
SASA = solvent-accessible surface area; ipTM = interface predicted TM score. Pose variation describes generated predictions, not a molecular-dynamics trajectory.
This inventory does not establish that every descriptor predicts experimental activity.}
\end{document}
